\documentclass[onsided]{book}
\usepackage[backend=biber,natbib=true,style=authoryear]{biblatex}
\addbibresource{/home/nguyen/1_NQBH/math/bib.bib}
\usepackage[utf8]{inputenc}
\usepackage{float}
\usepackage{graphicx}
\usepackage[colorlinks=true,linkcolor=blue,urlcolor=red,citecolor=magenta]{hyperref}
\usepackage{amsmath,amssymb,amsthm}
\allowdisplaybreaks
\numberwithin{equation}{section}
\newtheorem{assumption}{Assumption}[section]
\newtheorem{lemma}{Lemma}[section]
\newtheorem{corollary}{Corollary}[section]
\newtheorem{definition}{Definition}[section]
\newtheorem{proposition}{Proposition}[section]
\newtheorem{theorem}{Theorem}[section]
\newtheorem{notation}{Notation}[section]
\newtheorem{remark}{Remark}[section]
\newtheorem{example}{Example}[section]
\newtheorem{ques}{Question}[section]
\newtheorem{problem}{Problem}[section]
\newtheorem{conjecture}{Conjecture}[section]
\usepackage[left=0.5in,right=0.5in,top=1.5cm,bottom=1.5cm]{geometry}
\usepackage{fancyhdr}
\pagestyle{fancy}
\fancyhf{}
\addtolength{\headheight}{0pt}% obsolete
\lhead{\itshape \scriptsize \chaptername~\thechapter}
\rhead{\itshape \scriptsize \nouppercase{\leftmark}} %\nouppercase !
\renewcommand{\chaptermark}[1]{\markboth{#1}{}}
\cfoot{\thepage}

\title{Marie Sk\l odowska-Curie Actions}
\author{Collector: Nguyen Quan Ba Hong}
\date{\today}

\begin{document}
\maketitle
\setcounter{secnumdepth}{6}
\setcounter{tocdepth}{6}
\tableofcontents

%------------------------------------------------------------------------------%

\part{Wikipedia's}

\chapter{Wikipedia/Marie Sk\l odowska-Curie Actions}

\textit{Marie Skłodowska-Curie Actions (MSCA)} are a set of major research fellowships created by the \href{https://en.wikipedia.org/wiki/European_Union}{European Union}/\href{https://en.wikipedia.org/wiki/European_Commission}{European Commission} to support research in the \href{https://en.wikipedia.org/wiki/European_Research_Area}{European Research Area} (ERA).

The Marie Skłodowska-Curie Actions are among Europe's most competitive and prestigious research and innovation fellowships.[1][2]

%
Established in 1996 as Marie Curie Actions and known since 2014 as Marie Skłodowska-Curie Actions, they aim to foster the career development and further training of researchers at all career stages.

The Marie Skłodowska-Curie Actions promote interdisciplinary research and international collaborations, supporting scientists from not only within Europe but also across the globe.

%
Marie Skłodowska-Curie Actions are currently financed through the 8th \href{https://en.wikipedia.org/wiki/Framework_Programmes_for_Research_and_Technological_Development}{Framework Programme for Research and Technological Development} (called \href{https://en.wikipedia.org/wiki/Horizon_2020}{Horizon 2020}) and belong to the so-called `first pillar' of Horizon2020: ``Excellent Science.''

Through this funding scheme, the \href{https://en.wikipedia.org/wiki/Research_Executive_Agency}{Research Executive Agency} (REA) has devoted over €6 billion to the Marie Skłodowska-Curie Actions between 2014 and 2020.

%
Since the launch of the programme in 1996, over 100,000 researchers had received MSCA grants by Mar 2017.[3][4]

To mark this milestone, the European Commission selected thirty highly-promising researchers (who achieved the highest evaluation scores in 2016)[3] to showcase the EU's actions dedicated to excellence and worldwide mobility in research.[5]

%
\href{https://en.wikipedia.org/wiki/Marie_Curie}{Marie Skłodowska-Curie} is the \href{https://en.wikipedia.org/wiki/Poles_in_France}{Polish-French} namesake of the programme and was the \href{https://en.wikipedia.org/wiki/List_of_female_Nobel_laureates}{1st female Nobel prize winner}.

The only person \href{https://en.wikipedia.org/wiki/Nobel_Prize#Multiple_laureates}{to win a Nobel Prize for contributions in 2 different sciences} (\href{https://en.wikipedia.org/wiki/Nobel_Prize_in_Physics}{physics} and \href{https://en.wikipedia.org/wiki/Nobel_Prize_in_Chemistry}{chemistry}), she was also the 1st - and only woman - to have been awarded a Nobel Prize twice.

\section{Types of funding}
Fellowships are awarded by the European Commission across scientific disciplines within the framework of Horizon 2020.

%
MSCA are grouped into the following schemes:
\begin{itemize}
    \item Research Networks (ITN),
    \item Individual Fellowships (IF),
    \item Research and Innovation Staff Exchanges (RISE),
    \item Co-funding of regional, national and international programs involving mobility (COFUND),
    \item European Researchers' Night (NIGHT).
\end{itemize}
Within the framework of Horizon 2020, which runs from 2014 through 2020, MSCA will award €6.16 billion in funding.

\section{External links}
\begin{itemize}
    \item \href{http://ec.europa.eu/research/mariecurieactions/}{Main Marie Sk\l odowska-Curie Actions website}
    \item \href{http://cordis.europa.eu/fp7/people/}{People programme (FP7)}
    \item \href{https://www.mariecuriealumni.eu/}{Marie Curie Alumni Association}\hfill$\square$
\end{itemize}

%------------------------------------------------------------------------------%

\chapter{\href{https://ec.europa.eu/research/mariecurieactions/actions/individual-fellowships_en}{European Commission/Marie Skłodowska-Curie Actions}}

\section{\href{https://ec.europa.eu/research/mariecurieactions/}{Home}}

\subsection{Marie Skłodowska-Curie Actions - Research Fellowship Programme}
The Marie Skłodowska-Curie actions support researchers at all stages of their careers, regardless of age and nationality.

Researchers working across all disciplines are eligible for funding.

The MSCA also support cooperation between industry and academia and innovative training to enhance employability and career development.

\section{\href{https://ec.europa.eu/research/mariecurieactions/msca-actions_en}{Actions}}

\subsection{MSCA Actions}
\begin{quotation}
    Along with individual fellowships, the MSCA help develop training networks, promote staff exchanges and fund mobility programmes with an international flavor.
\end{quotation}
Grants provided by Marie Skłodowska-Curie Actions are available for all stages of a researcher's career.

Fellows include PhD candidates and those carrying out more advanced research.

%
Because they encourage individuals to work in other countries, the MSCA make the whole world a learning environment.

They encourage collaboration and sharing of ideas between different industrial sectors and research disciplines - all to the benefit of the wider European economy.

MSCA also back initiatives that break down barriers between academia, industry and business.

In addition, they reach out to the public with events that promote the value - and fun side - of science.

%
These inter-related goals are reflected in the various types of MSCA actions.

%
It all adds up to a huge investment.

In fact, the EU has set aside EUR 6.16 billion to be spent by 2020 on researcher training and career development.

\textit{Choose your action:}

\subsubsection{Innovative Training Networks (\href{https://ec.europa.eu/research/mariecurieactions/actions/research-networks_en}{ITN})}
ITN drive scientific excellence and innovation.

They bring together universities, research institutes and other sectors from across the world to train researchers to doctorate level.

\subsubsection{Individual Fellowships (\href{https://ec.europa.eu/research/mariecurieactions/actions/individual-fellowships_en}{IF})}
IF are great opportunities if you are an experienced researcher looking to give your career a boost by working abroad.

They offer exciting new learning opportunities and a chance to add some sparkle to your CV.

\subsubsection{Research and Innovation Staff Exchange (\href{https://ec.europa.eu/research/mariecurieactions/actions/staff-exchange_en}{RISE})}
RISE funds short-term exchanges of personnel between academic, industrial and commercial organizations throughout the world.

%
It helps people develop their knowledge, skills and careers, while building links between organizations working in different sectors of the economy, including universities, research institutes and SMEs.

\subsubsection{Co-funding of regional, national and international programmes (\href{https://ec.europa.eu/research/mariecurieactions/actions/co-funding-programmes_en}{COFUND})}
COFUND provides organizations with additional financial support for their own researcher training and career development programmes.

%
The extra funds are available for new or existing schemes for training researchers abroad and across various sectors.

\subsubsection{European Researchers' Night (\href{https://ec.europa.eu/research/mariecurieactions/actions/european-researchers-night_en}{NIGHT})}
This is a Europe-wide public event dedicated to popular science and fun learning. It takes place each year in September. Around 30 countries and over 300 cities are involved.

\subsection{\href{https://ec.europa.eu/research/mariecurieactions/actions/research-networks_en}{Research Networks}}
\begin{center}
    Innovative Training Networks (ITN) drive scientific excellence and innovation.
\end{center}

\begin{flushright}
    Funding scheme: \textbf{MSCA-ITN}
\end{flushright}
Innovative Training Networks (ITN) drive scientific excellence and innovation.

They bring together universities, research institutes and other sectors from across the world to train researchers to doctorate level.

\subsubsection{Types of ITN}

\subsubsection{What can be funded?}

\subsubsection{What does the funding cover?}

\subsubsection{How do we apply?}

\section{How to$\ldots$}

\section{News \& Events}

\section{Resources}

\section{Contact}


%------------------------------------------------------------------------------%

\printbibliography[heading=bibintoc]
\end{document}